\newpage
\phantomsection
\label{prf:boolean-ring-commutativity-from-idempotence}
\LRAProofFor{prop:boolean-ring-commutativity-from-idempotence}

\begin{remark*}[Return]
\hyperref[prop:boolean-ring-commutativity-from-idempotence]{Return to Theorem}
\end{remark*}

\begin{proposition*}[Commutativity from Idempotence]
In any Boolean ring $R$,
\[
  pq=qp \quad \text{for all } p,q\in R.
\]
Hence every Boolean ring is a commutative ring.
\end{proposition*}

\begin{proof}[Professional Standard Proof]
\LRAProofBodyStart
The proof of Proposition~\ref{prop:boolean-ring-self-additive-inverse}
established
\[
  pq+qp=0 \qquad \text{for all } p,q\in R.
\]
Hence $qp=-(pq)$. By Proposition~\ref{prop:boolean-ring-self-negation},
applied to the element $pq$, we have $-(pq)=pq$. Therefore $qp=pq$.
\end{proof}

\begin{proof}[Detailed Learning Proof]
\LRAProofBodyStart
The characteristic-$2$ proof produced a stronger intermediate identity:
$pq+qp=0$. This says that $qp$ is the additive inverse of $pq$. But in a
Boolean ring every element is its own negative, so the additive inverse of
$pq$ is just $pq$ itself. Therefore $qp=pq$.
\end{proof}

\begin{remark*}[Proof structure]
Reuse the intermediate identity from the characteristic-$2$ proof, then
replace additive inverse by self-negation.
\end{remark*}

\begin{dependencies}
\begin{itemize}
  \item \hyperref[prop:boolean-ring-self-additive-inverse]{Self Additive Inverse / Characteristic 2}.
  \item \hyperref[prop:boolean-ring-self-negation]{Self Negation}.
  \item \hyperref[def:notes-rings-definition-l65]{Commutative Ring}.
\end{itemize}
\end{dependencies}

\clearpage
